\section{Three passages}
\label{sec:cases}

Every quotation below is word for word from the released corrected transcript, with its punctuation and standard capitalization, and was checked by the author against the audio.

\subsection{The corpus matched a participant's transcript: \conv{871-004}, March 7, 1973}
\label{sec:case-dean}

With L. Patrick Gray's confirmation as FBI director going badly \citep{crs2018fbi}, Nixon told John Dean he wanted friendly senators to question Gray under oath about what J. Edgar Hoover had told him. Dean printed the passage from his own transcription \citep[Conversation No.~871-4]{dean2014defense}:

\begin{quote}\small
\textbf{Dean (2014):} \enquote{I'm going to make him lie, because I think Gray's not handling himself well. What he said yesterday with regard to the fact that this `jolly well' bullshit \ldots\ that it was highly improper for you to be present \ldots}

\textbf{This corpus (5:26--5:48):} \machine{I told him that Hoover had told me, and I'm going to make him lie because I think Gray's not handling himself well. What he said yesterday with regard to the fact that, it was, that, you know, this jolly well bullshit and all that sort of thing that you were, that it was highly improper for you to be present.}
\end{quote}

These are the recognizers' own words: the correction step made no edit here. Over the surrounding minutes the merged layer shared 95 of its 1,288 five-word sequences with Dean's text, against 10 to 20 for unrelated stretches. Dean printed what he selected; the corpus holds the whole conversation, including Nixon's question at 4:55 about the FBI's 1968 \machine{electronic surveillance of Vice President Agnew is phone when he was in Phoenix}. The documented check concerned calls from Albuquerque \citep{johnson1968deloach,hughes2014chasing}: the recognizer heard Nixon correctly, and the slip was his.

\subsection{The cover story in Nixon's words: \conv{916-019}, May 11, 1973}
\label{sec:case-916}

Eleven days after Haldeman and John Ehrlichman resigned, and days after Vernon Walters testified that they had sought the CIA's help in stopping an FBI investigation \citep{brinkley2015nixontapes1973}, Nixon met Haldeman and Alexander Haig in the Oval Office. At 10:20 Nixon (confirmed by the author's listening) told Haldeman how Ehrlichman's account of the June 23, 1972 meeting with the CIA's leaders should go:

\begin{quote}\small
\machine{Now, my point is, John's story should be perfect, should be, and I hope you talked to John. John's story is perfect. John should say his only concern was the whole CIA involvement and so forth and wanted to be sure that the, this investigation of Watergate didn't get out, as he put it, didn't get out of hand. [\ldots] Getting out of hand had to do with exactly what's now happening. Yeah. Which is unloading all this Pentagon papers. Investigation type stuff, the leak, the plumbers operation type stuff, the, Exactly. Wiretapping of the reporters and, and staff members type stuff.}
\end{quote}

\enquote{John} is Ehrlichman: the conversation placed \machine{Ehrlichman and me} at the meeting and said \machine{Dean was not there} and that \machine{Ehrlichman did do all the talking}. Parakeet transcribed these words where Whisper wrote nothing; the correction step did not touch them. At 13:00 Nixon named his concern and Haldeman answered (\machine{paint a} is the author's correction by ear of \machine{faint}):

\begin{quote}\small
\textbf{Nixon:} \machine{I, I don't need to paint a picture. I just want to be sure that John doesn't indicate that the President directed him to get Walters to get Gray to turn this off and set up a phony story}\\[2pt]
\textbf{Haldeman:} \machine{John would never indicate that. [\ldots] What the President tells us to do is not, not public, public consumption.}
\end{quote}

It was Nixon who, on June 23, 1972, told Haldeman that the CIA's leaders \enquote{should call the FBI in and say that we wish for the country, don't go any further into this case, period!} \citep{wspf1974exhibit1}. No published transcript of this exchange\footnote{Easy Nixon's unreviewed machine transcript of \conv{916-019} \citep{fine2026easynixon} does not contain it.} was found in the registry, in either Brinkley and Nichter volume, in \citet{dean2014defense} or on the Nixon Library's and Miller Center's pages for the conversation; \citet{kutler1997abuse} and the Presidential Recordings Digital Edition were not checked.

\subsection{The ordinary record, and further leads}
\label{sec:case-ordinary}

Most of the tapes are routine, and the corpus captures that too. On March 14, 1972, after John Mitchell's ITT testimony, Nixon telephoned Charles Colson and tried out a line, \machine{this is a smear a day Congress rather than a law day of Congress}; Colson laughed, and they agreed to use it (\conv{021-068}). No subject log records such exchanges. Further leads were surfaced by ranking corpus windows outside the registry for politically charged terms; each was heard and confirmed by the author (Table~\ref{tab:leads}).

\begin{table}[t]
\centering
\caption{Further leads, with prior-publication checks. Three other candidates were already in print \citep{dean2001rehnquist,dean2014defense}.}
\label{tab:leads}
\small
\begin{tabularx}{\linewidth}{@{}l l X X@{}}
\toprule
Conversation & Date, place & Passage & Prior publication \\
\midrule
\conv{789-006} & Sept.~30, 1972, Oval Office & \machine{super apathy among the blacks}, called \machine{just exactly where you want it} before the election & Not in the registry or \citet{brinkley2014nixontapes} \\
\conv{233-010} & Dec.~6, 1972, Camp David & \machine{make Sweden the black post. I want a black for Sweden} & Not in \citet{brinkley2014nixontapes}; in earlier machine transcripts \citep{fine2026easynixon} \\
\conv{690-011} & Mar.~21, 1972, Oval Office & An ambassadorship and a request to \machine{get at least 150 out of it} & Not in \citet{brinkley2014nixontapes}; in earlier machine transcripts \citep{fine2026easynixon} \\
\bottomrule
\end{tabularx}
\end{table}
